% theory.tex -- Theory subsections under Materials and Methods
\subsection{From Hopfield Energy to Multiplicity-Weighted Score Function}\label{sec:theory}

Let $\vX = [\vm_1, \ldots, \vm_K] \in \R^{d \times K}$ be a memory matrix whose columns are
unit-norm, PCA-encoded protein sequences from a family alignment of $K$ members, each represented
in $d$ dimensions. The modern Hopfield energy~\cite{ramsauerHopfieldNetworksAll2021} at state
$\vxi \in \R^d$ is
\[
E(\vxi) = \tfrac{1}{2}\|\vxi\|^{2}
- \tfrac{1}{\beta}\log\sum_{k}\exp(\beta\,\vm_{k}^{\!\top}\vxi),
\]
where $\beta > 0$ is an inverse temperature. Its score function,
\[
\nabla_{\vxi}\log p_{\beta}(\vxi)
= \beta[\vX\,\softmax(\beta\,\vX^{\!\top}\vxi) - \vxi],
\]
is exact and requires one attention operation. Applying the unadjusted Langevin algorithm
(ULA) with score step size $\alpha/\beta$, where $\alpha \in (0,1)$, gives the stochastic
attention update~\cite{varnerSAProtein2026}:
\begin{equation}\label{eq:ula}
  \vxi_{t+1}
    = (1-\alpha)\,\vxi_{t}
    + \alpha\,\vX\,\softmax\!\bigl(\beta\,\vX^{\!\top}\vxi_{t}\bigr)
    + \sqrt{2\alpha/\beta}\;\veps_{t},
  \qquad \veps_{t} \sim \mathcal{N}(\mathbf{0},\,\mathbf{I}_{d}).
\end{equation}
We extend this framework to condition generation on a functional subset by assigning a
multiplicity weight $r_k > 0$ to each stored pattern. The \emph{multiplicity-weighted} Hopfield
energy is
\begin{equation}\label{eq:weighted-energy}
  \boxed{
  E_{\vr}(\vxi) = \tfrac{1}{2}\|\vxi\|^{2}
    - \tfrac{1}{\beta}\log\sum_{k=1}^{K} r_k\,\exp\!\bigl(\beta\,\vm_{k}^{\!\top}\vxi\bigr),}
\end{equation}
where $\vr = (r_1, \ldots, r_K)^{\!\top}$.
When $r_k = 1$ for all $k$, Eq.~\eqref{eq:weighted-energy} reduces to the standard energy, and
patterns with higher multiplicity deepen the energy wells in their vicinity. For integer $r_k$
this is exactly equivalent to storing $r_k$ copies of pattern $\vm_k$ without the
$\mathcal{O}(d \cdot \sum_k r_k)$ memory cost, and for positive real $r_k$ it is the continuous
generalization of that replication.

\begin{proposition}[Exact Gaussian-mixture target]\label{prop:gmm-target}
For unit-norm memories, the Gibbs density associated with
Eq.~\eqref{eq:weighted-energy}, and hence the stationary law of the corresponding
continuous-time overdamped Langevin diffusion, is
\begin{equation}\label{eq:gmm-target}
  \boxed{
  p_{\vr}(\vxi)
  = \sum_{k=1}^{K} w_k\,
    \mathcal{N}\!\left(\vxi;\vm_k,\beta^{-1}\mathbf{I}_d\right),
  \qquad
  w_k = \frac{r_k}{\sum_j r_j}.}
\end{equation}
Independent equilibrium draws therefore require only
$k\sim\operatorname{Categorical}(\mathbf{w})$ followed by
$\vxi=\vm_k+\beta^{-1/2}\veps$, with
$\veps\sim\mathcal{N}(\mathbf{0},\mathbf{I}_d)$.
\end{proposition}

The derivation is given in the Supporting Information.
For unit-norm memories we can therefore draw equilibrium samples directly from the
mixture, and we use those draws as the reference in the exact-equilibrium benchmark.
We still run ULA for the reported generation runs, not because the target requires
it, but to stay comparable with the original stochastic-attention workflow and to
measure what finite steps cost.
Equation~\eqref{eq:gmm-target} is exact for the idealized dynamics that
move smoothly in time. The updates we actually run, Eqs.~\eqref{eq:ula} and
\eqref{eq:weighted-ula} below, advance in finite jumps of size $\alpha$, and a
stepped chain does not in general settle on the same distribution as the smooth one.
Two errors therefore remain at nonzero $\alpha$, a step-size bias that persists no
matter how long we run, and the error from stopping after finitely many
steps~\cite{durmusMoulinesULA2017}.

\begin{proposition}[Score function]\label{prop:score}
The score function of $p_{\vr}(\vxi) \propto \exp(-\beta\,E_{\vr}(\vxi))$ is
\[
\nabla_{\vxi}\log p_{\vr}(\vxi)
= \beta[\vX\,\softmax(\beta\,\vX^{\!\top}\vxi + \log\vr) - \vxi],
\]
where $\log\vr = (\log r_1, \ldots, \log r_K)^{\!\top}$ is the elementwise log-multiplicity.
\end{proposition}

The proof follows from the identity $r_k\,\exp(a_k) = \exp(a_k + \log r_k)$: the weights are
absorbed into the softmax logits. The resulting Langevin update is
\begin{equation}\label{eq:weighted-ula}
  \boxed{
  \vxi_{t+1}
    = (1-\alpha)\,\vxi_{t}
    + \alpha\,\vX\,\softmax\!\bigl(\beta\,\vX^{\!\top}\vxi_{t} + \log\vr\bigr)
    + \sqrt{2\alpha/\beta}\;\veps_{t}.}
\end{equation}
The \emph{only} change from Eq.~\eqref{eq:ula} is the addition of $\log r_k$ to each logit
before the softmax. The memory matrix $\vX$ is unchanged, the noise schedule is unchanged, and
the sampler retains its $\mathcal{O}(dK)$ per-step cost.

\subsection{Multiplicity Ratio and the Entropy Crossover}\label{sec:phase}

The user designates $K_{\mathrm{des}}$ of the $K$ stored patterns, leaving the
remaining $K_{\mathrm{bg}} = K - K_{\mathrm{des}}$ as background. Each stored pattern is one
sequence of the alignment, so $K_{\mathrm{des}}$ counts sequences and not alignment positions.
In the Kunitz family studied here, the memory holds $K = 99$ aligned domains, of which the
$K_{\mathrm{des}} = 32$ carrying Lys or Arg at the P1 contact position are designated and the
remaining $K_{\mathrm{bg}} = 67$ are background. The designation is
an input to the method rather than something the method infers, and nothing below
depends on what it means. A list of accession IDs would serve as well as a marker residue,
because the sampler sees only the resulting labels. We give every designated pattern the same
multiplicity and every background pattern the same multiplicity, and we write $\rho$
for the ratio between the two. Fixing the background multiplicity at $1$ sets the
scale, so $r_k = \rho$ for designated patterns and $r_k = 1$ for background
patterns, and $\rho$ is then the only free parameter.

The designated patterns therefore carry total multiplicity $K_{\mathrm{des}}\rho$
out of $K_{\mathrm{des}}\rho + K_{\mathrm{bg}}$ in all, so by
Eq.~\eqref{eq:gmm-target} the probability that a draw comes from a designated
component is
\begin{equation}\label{eq:f-eff}
  f_{\mathrm{eff}}(\rho)
  = \frac{K_{\mathrm{des}}\,\rho}{K_{\mathrm{des}}\,\rho + K_{\mathrm{bg}}} ,
\end{equation}
which we call the effective designated fraction. At $\rho=1$ it equals
$K_{\mathrm{des}}/K$, the designated share of the family, and as $\rho \to \infty$
it approaches $1$. Inverting Eq.~\eqref{eq:f-eff} gives the multiplicity ratio
needed for a target fraction $f$:
\begin{equation}\label{eq:rho-inverse}
  \rho(f) = \frac{f \cdot K_{\mathrm{bg}}}{K_{\mathrm{des}}\,(1-f)}.
\end{equation}
Choosing $\rho$ fixes the mixture we are aiming at, but it says nothing about how
sharply the sampler retrieves from that mixture. That is the job of $\beta$, and
$\beta$ has to be chosen before we can run anything. The division of labor is clean in
one direction: for unit-norm memories $\beta$ sets the width $\beta^{-1}$ shared by
every mixture component, so it controls how much neighboring components overlap, which
component a given state most likely came from, and therefore what the decoder returns,
but it does not touch the mixture weights $w_k$ or $f_{\mathrm{eff}}$. It is not clean
in the other direction, and the rest of this subsection is about how $\rho$ moves the
$\beta$ we should use.

The choice matters because both extremes of $\beta$ are useless. The attention weights
in Eq.~\eqref{eq:weighted-ula} are $\softmax(\beta\vX^{\!\top}\vxi + \log\vr)$. When
$\beta$ is small the logits are nearly flat, every stored pattern draws close to its
multiplicity weight $w_k$, and the sampler steps toward a weighted average of the whole
family, which decodes to a consensus-like blur. When $\beta$ is large the softmax
collapses onto whichever pattern the state is closest to, and the sampler sits on that
stored sequence and reproduces it. Useful generation happens between these two
regimes, and the attention entropy is what tells us where the transition is.

Following the base SA method~\cite{varnerSAProtein2026}, we track the Shannon entropy
$H_{\vr}(\beta)$ of the attention weights as $\beta$ increases, and we operate at the
onset of its drop. We call that drop the entropy crossover, and we locate its onset at
the point of maximum downward curvature of $H_{\vr}$ against $\log\beta$, writing the
inverse temperature there as $\beta^{*}$ (Fig.~\ref{fig:phase-transition}). This is a
convention for picking a reproducible operating point, not a claim that anything
thermodynamic happens there.

For standard SA the operating point does not have to be found by sweeping at all. The
companion paper~\cite{varnerSAProtein2026} establishes, from a concentration-of-measure
argument on the attention entropy, that the unweighted sampler ($\rho=1$) satisfies
\begin{equation}\label{eq:beta-star}
  \beta^{*} \;\approx\; 1.57 + 0.28\sqrt{d},
\end{equation}
with $R^{2}=0.97$ across eight Pfam families spanning $d \in [18, 186]$, and nearly
independent of $K$ for $K \geq 30$. The alignment alone therefore fixes $\beta^{*}$.

Multiplicity weighting breaks that shortcut, which is why $\rho$ and $\beta^{*}$ belong
in the same subsection. The biases $\log r_k$ sit in the logits before $\beta$ does
anything, so at $\beta=0$ the attention is already uneven: its weights are
$w_k = r_k/\sum_j r_j$ and its entropy is $H_{\vr}(0) = -\sum_k w_k \log w_k$, below
the uniform value $\log K$ that standard SA starts from. Every entropy curve therefore
starts lower and its drop moves to larger $\beta$
(Fig.~\ref{fig:phase-transition}), so for $\rho > 1$ we abandon
Eq.~\eqref{eq:beta-star} and re-find $\beta^{*}(\rho)$ by sweeping $\beta$ at each
$\rho$. Every experiment below runs at the $\beta^{*}(\rho)$ found this way. On the
Kunitz family it rose monotonically from 4.4 at $\rho=1$ to 9.3 at $\rho=1{,}000$.

Alongside $\beta^{*}$ we report the effective pattern count
$K_{\mathrm{eff}}(\vr) = (\sum_k r_k)^{2}/\sum_k r_k^{2}$, which equals $K$ when the
multiplicities are all equal and falls toward $K_{\mathrm{des}}$ as $\rho\to\infty$.
On Kunitz it fell from 99 to 32.1 over the range above. It summarizes how concentrated
the weights have become, and we report it beside the $\beta^{*}$ shift as a companion
observation rather than derive either from the other. The two are not the same
quantity, since $\log K_{\mathrm{eff}}$ is the Renyi-2 entropy of the weights and
equals the Shannon entropy $H_{\vr}(0)$ only when the weights are equal.

\subsection{The Calibration Gap}\label{sec:gap}

The signal chain from multiplicity weights to a decoded sequence marker passes through several
transformations:
\begin{equation}\label{eq:signal-chain}
  \rho
  \;\xrightarrow{\;\text{exact}\;}\;
  f_{\mathrm{eff}}
  \;\xrightarrow{\;\text{exact}\;}\;
  \bar{a}_{\mathrm{des}}
  \;\xrightarrow{\;\text{lossy}\;}\;
  f_{\mathrm{soft}}
  \;\xrightarrow{\;\text{nonlinear}\;}\;
  f_{\mathrm{obs}},
\end{equation}
where $\bar{a}_{\mathrm{des}}$ is the mean attention weight on designated patterns during
sampling, $f_{\mathrm{soft}}$ is the continuous-valued residue probability at the marker position,
and $f_{\mathrm{obs}}$ is the fraction of generated sequences displaying the designated marker
after argmax decoding. We define the calibration gap $\Delta = f_{\mathrm{eff}} - f_{\mathrm{obs}}$
and decompose it into three independently measurable components:
\begin{equation}\label{eq:gap-decomp}
  \Delta = \underbrace{(f_{\mathrm{eff}} - \bar{a}_{\mathrm{des}})}_{\Delta_{\mathrm{attn}}\;\approx\;0}
  + \underbrace{(\bar{a}_{\mathrm{des}} - f_{\mathrm{soft}})}_{\Delta_{\mathrm{PCA}}}
  + \underbrace{(f_{\mathrm{soft}} - f_{\mathrm{obs}})}_{\Delta_{\mathrm{argmax}}}.
\end{equation}
Because the target is exactly a mixture, we can also describe the final step by
conditioning on which component actually generated the sample. Let $C$ be that
component and let $\mathcal{D}$ be the index set of the designated patterns, so that
$C\in\mathcal{D}$ says the draw came from a designated component. Define
\[
\begin{aligned}
q_{\mathrm{des}}
  &= \Pr(\text{decoded marker positive}\mid C\in\mathcal{D}),\\
q_{\mathrm{bg}}
  &= \Pr(\text{decoded marker positive}\mid C\notin\mathcal{D}).
\end{aligned}
\]
Then
\begin{equation}\label{eq:decoder-pushforward}
  f_{\mathrm{obs}}
  = f_{\mathrm{eff}}q_{\mathrm{des}}
  + (1-f_{\mathrm{eff}})q_{\mathrm{bg}},
  \qquad
  \Delta
  = f_{\mathrm{eff}}(1-q_{\mathrm{des}})
  - (1-f_{\mathrm{eff}})q_{\mathrm{bg}}.
\end{equation}
Thus $f_{\mathrm{obs}}=f_{\mathrm{eff}}$ is not something the algebra requires at
$\rho=1$. It holds only when the decoder's false negatives happen to balance its
false positives. The calibration gap is therefore a quantity to be measured across
the whole path from latent state to sequence, not a diagnostic for tuning that
path, and matching $f_{\mathrm{obs}}$ to $f_{\mathrm{eff}}$ in particular does not
identify a correct PCA scaling, because a decoder that carries no information at
all can match the two by trading one error rate against the other.
To describe the PCA geometry across families, we define the separation index
$S = (\bar{c}_{\mathrm{within}} - \bar{c}_{\mathrm{between}}) /
[\tfrac{1}{2}(\sigma_{\mathrm{within}} + \sigma_{\mathrm{between}})]$, where $\bar{c}$ and
$\sigma$ denote the mean and standard deviation of pairwise cosine similarities within and
between functional groups in PCA space. This is a Fisher-type ratio defined for this study
on pairwise cosine similarities, not the classical Fisher discriminant criterion,
which is computed from class means and scatter matrices.
